\chapter{Methodology}
\label{ch:methodology}

\section{Research Design}

This research follows a controlled experimental design, consistent with Type 5 (Experimental) dissertations under the Department's classification. A fixed dataset of 300 synthetic user profiles and a curated catalogue of 20 real UK credit cards form the basis for all experiments. The same 300 profiles are evaluated under three prompting conditions, and outputs are assessed against a content-based recommender baseline. The evaluation combines quantitative scoring (recommendation overlap, top-1 accuracy, deterministic factual correctness) with an automated qualitative criterion (explanation quality, scored via LLM-as-judge --- Section~\ref{sec:evalframework}), producing a mixed-method analysis of LLM recommendation quality. The study is entirely computational: it does not involve human participants or real personal financial data (Section~\ref{sec:ethics}).

\section{Dataset Construction}

\subsection{Credit Card Catalogue}

A structured catalogue of 20 UK credit cards was assembled from publicly available issuer sources (see \texttt{docs/card-sourcing-notes.md} for provenance). Each entry records the card name, issuer, network, annual fee, reward type, reward rate by spending category (groceries, fuel, dining, travel, online shopping), welcome offer value, foreign transaction fee, and APR range. This catalogue defines the fixed product space from which all recommendations --- both the baseline's and the LLM's --- are drawn; no card outside this set can be a valid recommendation, which is itself checked as part of the evaluation pipeline (Section~\ref{sec:evalframework}).

\subsection{Synthetic User Profiles}

All user profiles are synthetically generated; no real personal financial data is used at any stage (Section~\ref{sec:ethics}). An initial 80-profile dataset, using hand-specified archetype ranges, was expanded to 300 profiles following supervisor feedback that a larger, more rigorously grounded sample was needed for adequate statistical power in the planned significance testing.

Profiles are built around eight fixed archetypes representing distinct UK credit-card consumer segments: student (low spend), family (grocery-heavy), frequent traveller, balanced cashback user, large-purchase planner, online shopper, fuel commuter, and dining-heavy user. Rather than relying solely on hand-specified ranges, spending distributions across six categories (groceries, travel, fuel, dining, online shopping, other) are calibrated against the ONS \textit{Living Costs and Food Survey 2022--23} (Office for National Statistics, 2023), so that each archetype's spending pattern is anchored to an empirically observed UK household expenditure distribution rather than an arbitrary one (see \texttt{docs/data-justification.md} for the specific ONS figures used per category). This is consistent with the wider literature on financial-recommendation research, where individual-level real spending data is not publicly available for use due to UK GDPR and Financial Services and Markets Act 2000 restrictions, and simulated or anonymised profiles are the standard methodological choice.

Profiles are distributed 38/38/38/38/37/37/37/37 across the eight archetypes (summing to exactly 300), generated deterministically via a fixed random seed (\texttt{seed = 42}) for full reproducibility. Each profile contains: a profile identifier, archetype label, age group, monthly income, monthly spend across six categories, annual fee preference, primary goal, travel interest level, spending style, preferred reward type, a cost-sensitivity flag, and a planned large-purchase amount (non-zero only for the large-purchase-planner archetype).

\section{Baseline Method}
\label{sec:baseline}

An initial baseline used four self-defined weighted heuristics (reward alignment, fee suitability, welcome-offer relevance, large-purchase suitability). Following supervisor feedback that this baseline was not grounded in established methodology, it was replaced with a content-based cosine-similarity approach, a well-established method in the recommender-systems literature (Lops, de Gemmis and Semeraro, 2011; Pazzani and Billsus, 2007; see Section~2.3).

Each credit card is represented as a numeric feature vector over eight attributes: annual fee (inverted, so a lower fee scores higher), base reward rate, category-specific reward rates (grocery, fuel, dining, travel, online shopping), and foreign transaction fee (inverted). Each user profile is represented as a preference vector over the same eight dimensions: an annual-fee-preference weighting derived from the profile's stated fee tolerance, and category weightings derived from the profile's monthly spend in each category (normalised to the profile's own maximum category spend). Both card and profile vectors are min--max scaled to $[0, 1]$. For each profile, cosine similarity is computed between the profile vector and every card vector in the catalogue, and the three highest-scoring cards constitute the baseline recommendation, recorded in \texttt{data/ground\_truth.csv}. This baseline is fully deterministic, transparent, reproducible, and independent of any language model, making it an appropriate conventional comparator for the LLM-generated recommendations evaluated in this study.

\section{LLM Experimental Setup}

\subsection{Model and Infrastructure}

The language model used to generate recommendations in this study is Meta's Llama 3.1 (8B-parameter, instruction-tuned variant), accessed through the Groq inference API on its free tier. Groq was selected for its low-latency inference at no financial cost, appropriate for an academic project with resource constraints; Groq's LPU (Language Processing Unit) infrastructure is purpose-built for LLM inference and offers sub-second time-to-first-token (Groq, 2024). The free tier used in this study is subject to a rate limit (requests per minute and tokens per minute) and a daily request cap, which materially affected the practical duration of the experimental runs (see Section~\ref{sec:limitations}). The model is used in standard inference mode with no fine-tuning or architectural modification; temperature was fixed at 0.2 across all calls to reduce (but not eliminate) run-to-run variance, with the residual variance itself the subject of the consistency evaluation (Section~\ref{sec:consistency-rationale}).

\subsection{Prompting Strategies}

Three prompt engineering strategies are evaluated across all 300 profiles:

\begin{itemize}
    \item \textbf{Zero-shot prompting}: the model receives the user profile and the full card catalogue and is instructed to recommend the top three cards, with one short justification per card, without any worked examples.
    \item \textbf{Structured prompting}: the model receives the same input alongside an explicit JSON output schema specifying card ID, card name, a one-sentence reason, and any relevant warnings for each of the three ranked recommendations.
    \item \textbf{Few-shot prompting}: the model is provided with two complete worked examples --- a profile alongside its correct recommendation and justification --- before being presented with the target profile, to enable in-context learning (Brown et al., 2020; see Section~2.2).
\end{itemize}

All 300 profiles are evaluated under each of the three strategies, producing 900 LLM inference calls in total. Each result is stored with its profile identifier, strategy label, up to three extracted recommended card IDs, the full raw model response, and inference latency.

\textbf{A methodological correction made during development:} the zero-shot and few-shot prompt templates originally instructed the model to return bare card IDs only, with explicit instructions not to provide any explanation. This was identified as inconsistent with the evaluation framework's ``explanation quality'' criterion (Section~\ref{sec:evalframework}), which requires reasoning text to score. Both prompts were revised to require a one-sentence justification per recommended card, matching the reasoning already present in the structured condition, and the zero-shot and few-shot experimental runs were repeated under the corrected prompts. This correction was completed successfully: all 900 responses (300 profiles $\times$ 3 strategies) were verified to contain zero duplicate profile--strategy pairs and, for zero-shot and few-shot, to be in the corrected reasoning-required format, before any scoring was performed.

\subsection{Evaluation Framework}
\label{sec:evalframework}

LLM outputs are assessed against four criteria:

\begin{enumerate}
    \item \textbf{Recommendation quality} --- the overlap between the LLM's top-3 recommendation and the baseline's top-3 recommendation for the same profile, reported as an overlap score (proportion of the three baseline cards also present in the LLM's top three) and as top-1 exact-match accuracy.
    \item \textbf{Factual correctness} --- accuracy of specific, checkable card details cited in the response (card name, stated annual fee, stated reward-rate percentages), checked deterministically against \texttt{data/cards.csv}. Unlike the original methodology's framing of this as a criterion requiring human judgement, a deterministic checker was implemented instead (\texttt{src/factual\_correctness.py}): it extracts specific numeric/named claims from each response and compares them directly against ground truth, which is more reproducible and removes rater subjectivity for the subset of claims that are objectively checkable (see Section~\ref{sec:limitations} for the scope and limits of this approach).
    \item \textbf{Explanation quality} --- clarity and apparent grounding of the model's stated reasoning, scored via an LLM-as-judge approach (Zheng et al., 2023; Section~2.4) rather than the originally planned single-rater human scoring, given the scale of the dataset (900 responses) and the documented reliability of strong-model judgement against human preference (Zheng et al., 2023 report over 80\% agreement, comparable to human--human agreement). The judge model is Llama 3.3 70B --- deliberately larger than, and architecturally distinct in scale from, the Llama 3.1 8B model being evaluated, to reduce self-enhancement bias (Zheng et al., 2023). For each response, the judge scores each of the three recommendation lines and an overall explanation-quality score on a 1--3 ordinal scale (1 = poor: missing, generic, or irrelevant reasoning; 2 = adequate: understandable but vague or only partially grounded; 3 = good: clear and explicitly grounded in the user's stated profile or the card's actual attributes), returned as structured JSON for reliable parsing (\texttt{src/llm\_judge.py}, \texttt{prompts/judge\_explanation\_quality.txt}). This substitution of model judgement for human judgement is treated explicitly as a methodological choice, not an oversight, and is discussed as a limitation in Section~\ref{sec:limitations}.
    \item \textbf{Consistency} --- stability of recommendations across repeated runs of the same profile and strategy, measured as the mean pairwise Jaccard similarity of the top-3 card sets across three runs, on a stratified subsample of profiles drawn evenly across all eight archetypes (\texttt{src/consistency\_check.py}; see Section~\ref{sec:consistency-rationale} for the sampling rationale).
\end{enumerate}

Where the original methodology specified a 3-point ordinal rubric aggregated via Kruskal-Wallis testing across all four criteria, the two criteria with automated, continuous-valued scoring (recommendation quality, factual correctness) are analysed using their native continuous/proportion scales, with Kruskal-Wallis used specifically for the ordinal/count-based comparisons across the three prompt strategies, and paired Wilcoxon signed-rank tests used for pairwise strategy comparisons on a per-profile basis (matching the profile-level pairing inherent in the experimental design, i.e.\ every profile is evaluated under all three strategies).

\subsection{Consistency Sampling Rationale}
\label{sec:consistency-rationale}

Repeating all 300 profiles across all three strategies would require 2,700 additional inference calls beyond the primary 900-call experiment. Given the free-tier rate and daily-request constraints observed during the primary data collection (Section~\ref{sec:limitations}), a stratified subsample is used instead: 3 profiles per archetype (24 profiles total, drawn with a fixed random seed for reproducibility) are each queried 3 times per strategy (216 additional calls in total), and stability is reported as the mean pairwise Jaccard overlap of the resulting top-3 sets. This ensures every archetype is represented in the consistency analysis without requiring a full-N repeat, a scope decision consistent with common practice in LLM stability studies where full repetition is cost-prohibitive.

\section{Data Analysis}

Scores are aggregated by prompt strategy. The Kruskal-Wallis H-test is used to test for an overall difference in recommendation-quality scores across the three strategies (appropriate for comparing more than two independent groups on non-normally-distributed ordinal/proportion data). Paired Wilcoxon signed-rank tests are additionally reported for each pairwise strategy comparison on a per-profile basis, since every profile is evaluated under all three conditions (a repeated-measures design), which the paired test is better powered to detect than the omnibus test alone. Per-archetype breakdowns are also reported, to check whether prompt-strategy effects are uniform across the eight consumer segments or archetype-dependent.

\section{Ethical Considerations}
\label{sec:ethics}

This study does not involve human participants, real financial data, or personally identifiable information at any stage. All user profiles are entirely synthetic; the real UK consumer expenditure statistics (ONS, 2023) used for calibration inform only the statistical distribution of category-level spending values used to generate synthetic profiles, not any individual-level record. The study adheres to standard ethical guidelines applicable to computational and data-science research and does not require formal ethical approval beyond departmental good-practice standards.

\section{Limitations}
\label{sec:limitations}

\textit{Full discussion in Chapter~\ref{ch:discussion}; summarised here per the standard dissertation structure.} The use of synthetic profiles means findings may not fully generalise to real consumer behaviour, even with ONS-calibrated distributions, since calibration does not capture irregular expenditure, life events, or behavioural nuances a real individual would exhibit. The study evaluates a single LLM (Llama 3.1, 8B parameters) for recommendation generation, which limits cross-model generalisability; a larger or differently-trained model may perform materially differently. Explanation-quality scoring via LLM-as-judge, while grounded in the literature (Zheng et al., 2023) and designed to mitigate self-enhancement bias through the use of a distinct, larger judge model, still substitutes model judgement for human judgement and inherits other documented biases of the approach (position and verbosity bias). The consistency evaluation uses a stratified subsample rather than the full 300 profiles, for the practical reasons given in Section~\ref{sec:consistency-rationale}. The deterministic factual-correctness checker (Section~\ref{sec:evalframework}) can only verify claims that are specific and numeric or exactly named; vague or qualitative claims in a response are neither confirmed nor flagged, so the reported factual-error rate is a lower bound on true factual inaccuracy, not a complete audit. The free-tier API's rate and daily-request limits meant the primary 900-call run, and the correction rerun described in Section~4.2, each took multiple days of wall-clock time to complete, which is a practical constraint on this study's methodology rather than a property of the models or prompting strategies themselves.
